\newpage

\section{量子力学：低能中子--中子散射与自旋依赖势}
\label{sec:nn-scattering}

\begin{cheatbox}[考场速查卡]
\begin{tabular}{lll}
\textbf{题目给的条件/要求} & \textbf{你该用的工具} & \textbf{绝对不要做的事} \\
求微分截面 $\sigma(\theta)$ & Born 近似：$|f^{(1)}|^2$ & 不要试图用光学定理求角分布 \\
求总截面（Born 框架内） & $\int |f^{(1)}|^2 d\Omega$ & 不要把 $f^{(1)}(0)$ 代入光学定理 \\
已知 $\sigma_{\text{tot}}$，求 $\text{Im}[f(0)]$ & 光学定理 & 不要试图用 Born 近似反推 \\
分波法验证自洽性 & 光学定理 + 分波公式 & 不要把 Born 和分波混用 \\
题目说``用光学定理'' & 光学定理 & 不要绕道 Born 近似 \\
\end{tabular}

一句话总结：Born 近似是``建造'' 微分截面的工具；光学定理是``检验'' 总截面是否幺正的尺子。一级 Born 的结果只能留在 Born 的框架内用（积分得总截面），绝不能跨到光学定理那边去（因为光学定理需要精确的 $f(0)$，而一级 Born 给不出虚部）。
\end{cheatbox}

\begin{modelbox}[低能非极化中子--中子散射]
两个中子发生低能散射（$E \to 0$），相互作用势为
\[
V(r) = \begin{cases}
V_0 \, \vec{\sigma}_1 \cdot \vec{\sigma}_2, & r < a \\[4pt]
0, & r > a
\end{cases}
\]
其中 $\vec{\sigma}_1, \vec{\sigma}_2$ 为两中子的泡利矩阵，$V_0$ 为常数。入射中子与靶中子均为非极化。求散射总截面 $\sigma_t$。
\end{modelbox}

\begin{tipbox}[考点快速分析]
\begin{itemize}
    \item 自旋耦合：$\vec{\sigma}_1 \cdot \vec{\sigma}_2$ 的本征值由总自旋 $S$ 决定
    \item 全同费米子统计：总波函数反对称；s 波空间对称 $\Rightarrow$ 自旋单态
    \item 低能极限：$E \to 0$ 时仅 $l = 0$ 分波有贡献；三重态被统计与离心势垒双重压制
    \item 非极化平均：单态概率 $1/4$，三重态概率 $3/4$
    \item 全同粒子对称化：$f(\theta) \to f(\theta) \pm f(\pi - \theta)$
\end{itemize}
\end{tipbox}

\begin{derivationbox}[自旋耦合、统计与有效势]
总自旋 $\vec{S} = \vec{S}_1 + \vec{S}_2 = \frac{\hbar}{2}(\vec{\sigma}_1 + \vec{\sigma}_2)$。由
\[
\vec{S}^2 = \vec{S}_1^2 + \vec{S}_2^2 + 2\vec{S}_1 \cdot \vec{S}_2 = \frac{3}{4}\hbar^2 + \frac{3}{4}\hbar^2 + 2\vec{S}_1 \cdot \vec{S}_2 = \frac{3}{2}\hbar^2 + 2\vec{S}_1 \cdot \vec{S}_2,
\]
得
\[
\vec{S}_1 \cdot \vec{S}_2 = \frac{1}{2}\left( \vec{S}^2 - \frac{3}{2}\hbar^2 \right).
\]

再利用 $\vec{S}_1 = \frac{\hbar}{2}\vec{\sigma}_1$，$\vec{S}_2 = \frac{\hbar}{2}\vec{\sigma}_2$，有
\[
\vec{\sigma}_1 \cdot \vec{\sigma}_2 = \frac{4}{\hbar^2}\, \vec{S}_1 \cdot \vec{S}_2 = \frac{2}{\hbar^2}\left( \vec{S}^2 - \frac{3}{2}\hbar^2 \right) = \frac{2\vec{S}^2}{\hbar^2} - 3.
\]

对总自旋量子数 $s = 0$（单态，$\vec{S}^2 = 0(0+1)\hbar^2 = 0$）和 $s = 1$（三重态，$\vec{S}^2 = 1(1+1)\hbar^2 = 2\hbar^2$）：
\begin{align*}
\vec{\sigma}_1 \cdot \vec{\sigma}_2 \, \Phi_{00} &= -3 \, \Phi_{00}, \\
\vec{\sigma}_1 \cdot \vec{\sigma}_2 \, \Phi_{1m} &= +1 \, \Phi_{1m}.
\end{align*}

因此势可按自旋通道分开：
\begin{itemize}
    \item \textbf{单态}（$s = 0$）：$V_S(r) = -3V_0 \, \Theta(a - r)$，等效为吸引方势阱（设 $V_0 > 0$）
    \item \textbf{三重态}（$s = 1$）：$V_T(r) = +V_0 \, \Theta(a - r)$，等效为排斥方势垒
\end{itemize}
\end{derivationbox}

\begin{derivationbox}[为什么要转入质心系？——两体问题分解]
实验室系中写出的两体薛定谔方程：
\[
\left[ -\frac{\hbar^2}{2m}\nabla_1^2 - \frac{\hbar^2}{2m}\nabla_2^2 + V(|\vec{r}_1 - \vec{r}_2|) \right] \Psi(\vec{r}_1, \vec{r}_2) = E_{\text{tot}} \Psi(\vec{r}_1, \vec{r}_2).
\]

引入质心坐标和相对坐标：
\[
\vec{R} = \frac{\vec{r}_1 + \vec{r}_2}{2}, \qquad \vec{r} = \vec{r}_1 - \vec{r}_2.
\]

下面验证动能的分解。先求坐标变换的逆：
\[
\vec{r}_1 = \vec{R} + \frac{\vec{r}}{2}, \qquad \vec{r}_2 = \vec{R} - \frac{\vec{r}}{2}.
\]

对 $\vec{r}_1$ 求梯度时，$\vec{R}$ 和 $\vec{r}$ 都随 $\vec{r}_1$ 变：
\[
\nabla_1 = \frac{\partial}{\partial \vec{r}_1} = \frac{\partial \vec{R}}{\partial \vec{r}_1} \frac{\partial}{\partial \vec{R}} + \frac{\partial \vec{r}}{\partial \vec{r}_1} \frac{\partial}{\partial \vec{r}} = \frac{1}{2}\nabla_R + \nabla_r.
\]
同理
\[
\nabla_2 = \frac{1}{2}\nabla_R - \nabla_r.
\]

计算拉普拉斯算符（注意 $\nabla_R$ 与 $\nabla_r$ 对易，因为它们对不同的独立变量求导）：
\begin{align*}
\nabla_1^2 &= \left( \frac{1}{2}\nabla_R + \nabla_r \right)^2 = \frac{1}{4}\nabla_R^2 + \nabla_R \cdot \nabla_r + \nabla_r^2, \\
\nabla_2^2 &= \left( \frac{1}{2}\nabla_R - \nabla_r \right)^2 = \frac{1}{4}\nabla_R^2 - \nabla_R \cdot \nabla_r + \nabla_r^2.
\end{align*}

两式相加，交叉项 $\pm \nabla_R \cdot \nabla_r$ 精确抵消：
\[
\nabla_1^2 + \nabla_2^2 = \frac{1}{2}\nabla_R^2 + 2\nabla_r^2.
\]

代回动能：
\[
-\frac{\hbar^2}{2m}\nabla_1^2 - \frac{\hbar^2}{2m}\nabla_2^2 = -\frac{\hbar^2}{4m}\nabla_R^2 - \frac{\hbar^2}{m}\nabla_r^2.
\]

定义
\[
M = 2m \text{（总质量）}, \qquad \mu = \frac{m_1 m_2}{m_1 + m_2} = \frac{m}{2} \text{（约化质量）},
\]
则 $4m = 2M$，$m = 2\mu$，上式化为
\[
-\frac{\hbar^2}{2m}\nabla_1^2 - \frac{\hbar^2}{2m}\nabla_2^2 = -\frac{\hbar^2}{2M}\nabla_R^2 - \frac{\hbar^2}{2\mu}\nabla_r^2.
\]

用动量算符 $\vec{P} = -i\hbar\nabla_R$，$\vec{p} = -i\hbar\nabla_r$ 写出就是
\[
\frac{\vec{p}_1^2}{2m} + \frac{\vec{p}_2^2}{2m} = \frac{\vec{P}^2}{2M} + \frac{\vec{p}^2}{2\mu}.
\]

由于势只依赖于 $r = |\vec{r}_1 - \vec{r}_2|$，方程可分离变量：
\[
\Psi(\vec{r}_1, \vec{r}_2) = \Phi(\vec{R}) \, \psi(\vec{r}).
\]
\begin{itemize}
    \item \textbf{质心运动：}$-\frac{\hbar^2}{2M}\nabla_R^2 \Phi = E_{CM} \Phi$。没有外势，质心做自由运动，与散射无关。
    \item \textbf{相对运动：}$\left[ -\frac{\hbar^2}{2\mu}\nabla_r^2 + V(r) \right] \psi = E\psi$。这正是一个质量为 $\mu$ 的粒子在势 $V(r)$ 中的散射问题。
\end{itemize}

为什么这让全同性更明显？在质心系中总动量为零，两中子以相同速率反向运动，地位完全平等。交换 $1 \leftrightarrow 2$ 对应
\[
\vec{r} = \vec{r}_1 - \vec{r}_2 \xrightarrow{1\leftrightarrow 2} \vec{r}_2 - \vec{r}_1 = -\vec{r}.
\]

若 $\vec{r}$ 的球坐标为 $(r, \theta, \phi)$，则 $-\vec{r}$ 的坐标为 $(r, \pi - \theta, \phi + \pi)$（由 $\cos\theta' = -\cos\theta$ 得 $\theta' = \pi - \theta$；$x, y$ 分量反向得 $\phi' = \phi + \pi$）。

为什么只关心 $\theta \to \pi - \theta$ 而不管 $\phi$？入射平面波 $e^{ikz}$ 已经选出了一个特殊方向（$z$ 轴），所以散射振幅 $f(\theta)$ 确实依赖于极角 $\theta$（它描述的是偏离入射方向的程度）。但系统仍具有绕 $z$ 轴的旋转对称性（$L_z$ 守恒），因此 $f$ 不依赖于方位角 $\phi$。既然物理结果与 $\phi$ 无关，$\phi \to \phi + \pi$ 不会产生任何可观测效应，交换的等价操作只剩下
\[
\theta \longrightarrow \pi - \theta.
\]

实验室系里``谁入射谁静止'' 的区分在质心系中完全消失。
\end{derivationbox}

\begin{insightbox}[为什么说这是全同粒子散射？——实验语言 vs 物理现实]
题目中``入射中子'' 与``靶中子'' 的说法只是实验语言的方便描述：初始时刻一束中子飞向一个静止的中子靶。但物理上：
\begin{enumerate}
    \item 两者都是中子：质量、电荷（零）、自旋完全相同，没有任何内禀标签可以区分它们。
    \item 质心系的对称性：转入质心系后，两中子的动量始终相反、地位完全平等。哈密顿量 $H = \frac{\vec{p}_1^2}{2m} + \frac{\vec{p}_2^2}{2m} + V(|\vec{r}_1 - \vec{r}_2|)$ 在交换 $1 \leftrightarrow 2$ 下严格不变。
    \item 探测器无法分辨：实验上探测器捕获一个中子，你没有任何办法知道它是原来那束里的还是靶里的。因此必须把两条散射路径（粒子 1 飞向 $\theta$ vs 粒子 2 飞向 $\theta$）的振幅叠加，产生干涉。
\end{enumerate}

一句话：``入射'' 和``靶'' 只描述 $t \to -\infty$ 时的初始制备；一旦散射发生，两个中子就是全同的、不可区分的费米子，总波函数必须反对称。
\end{insightbox}

\subsubsection*{为什么三重态在 $E \to 0$ 被压制？}

全同费米子要求总波函数反对称。s 波（$l = 0$）空间波函数对称，只能配反对称的自旋单态。三重态要求空间反对称，最低允许的分波是 $l = 1$（p 波）。p 波在原点有离心势垒 $l(l+1)\hbar^2/2\mu r^2$，相移 $\delta_1 \sim k^3$（当 $k \to 0$），截面被 $k^4$ 压制。因此在 $E \to 0$ 极限下，只有单态的 s 波有贡献。

\begin{derivationbox}[为什么这道题必须用分波法？]
\subsubsection*{分波法的``触发条件''——与 Born 近似的对比}

这道题有两个特征，直接决定了分波法是唯一自然的选择：
\begin{itemize}
    \item \textbf{低能 $E \to 0$：}Born 近似要求入射粒子动能远大于势能扰动（或势足够弱），其适用条件是 $ka \gg 1$ 且势弱。这里 $k \to 0$，Born 近似的前提崩塌。
    \item \textbf{有限程方势阱：}势只在 $r < a$ 内非零，具有明确的作用范围。分波法的核心思想是把平面波按球谐函数展开，每个分波 $l$ 独立感受有效势 $V(r) + \frac{l(l+1)\hbar^2}{2\mu r^2}$。当 $ka \ll 1$ 时，只有 $l = 0$ 的 s 波能``挤进'' 势场区域，所有 $l \ge 1$ 被离心势垒压制。这使得问题从一整个偏微分方程降维为一个常微分方程。
\end{itemize}

考场识别标签：看到``$E \to 0$'' + ``球对称势'' + ``求总截面''，反射弧应该是：分波法 $\to$ 只算 $l = 0$ $\to$ 求相移 $\delta_0$ $\to$ 总截面 $4\pi\sin^2\delta_0/k^2$。
\end{derivationbox}

\begin{derivationbox}[离心势垒与低能压制]
\subsubsection*{离心势垒与低能压制——$\delta_l \sim k^{2l+1}$ 从哪来？}

对第 $l$ 分波，径向方程为
\[
\frac{d^2 u_l}{dr^2} + \left[ k^2 - \frac{l(l+1)}{r^2} - \frac{2\mu V(r)}{\hbar^2} \right] u_l = 0.
\]

在势外（$r > a$）且 $kr \ll 1$ 时，$k^2$ 可忽略，方程近似为
\[
\frac{d^2 u_l}{dr^2} - \frac{l(l+1)}{r^2} u_l \approx 0.
\]

设 $u_l \sim r^\alpha$，得 $\alpha(\alpha - 1) = l(l+1)$，即 $\alpha = l+1$（取正则解，$\alpha = -l$ 因 $u_l(0) = 0$ 被舍弃）。因此
\[
u_l(r) \sim r^{l+1} \quad \Longrightarrow \quad R_l(r) = \frac{u_l}{r} \sim r^l.
\]
物理图像：$l \ge 1$ 的分波在原点被强制为零（$R_l(0) = 0$），$l$ 越大零点越``深''，粒子越难进入势的作用范围。

相移的标度行为（Born 近似对相移的积分公式）：
\[
\sin\delta_l \approx -\frac{2\mu k}{\hbar^2} \int_0^\infty r^2 V(r) \bigl[j_l(kr)\bigr]^2 \, dr.
\]
当 $kr \ll 1$ 时，球 Bessel 函数的渐近行为是 $j_l(kr) \sim (kr)^l$。代入得
\[
\sin\delta_l \sim k \cdot k^{2l} = k^{2l+1}.
\]
对 $l = 0$：$\delta_0 \sim k$；对 $l = 1$：$\delta_1 \sim k^3$；对 $l = 2$：$\delta_2 \sim k^5$。

截面的标度行为：
\[
\sigma_l = \frac{4\pi}{k^2}(2l+1)\sin^2\delta_l \sim \frac{1}{k^2} \cdot (k^{2l+1})^2 = k^{4l}.
\]
对 $l = 0$：$\sigma_0 \sim k^0 = \text{常数}$；对 $l = 1$：$\sigma_1 \sim k^4$；对 $l = 2$：$\sigma_2 \sim k^8$。

结论：低能时 s 波（$l = 0$）截面趋于常数，所有 $l \ge 1$ 的分波被 $k^{4l}$ 强烈压制。三重态因统计要求必须取 $l \ge 1$，其最低分波 $l = 1$ 的截面 $\sim k^4 \to 0$，故在 $E \to 0$ 极限下无贡献。
\end{derivationbox}

\begin{derivationbox}[s 波分波法与相移]
\subsubsection*{径向方程从哪来？——逐步推导}

在质心系中，相对运动满足定态薛定谔方程（设势只依赖于 $r = |\vec{r}_1 - \vec{r}_2|$）：
\[
-\frac{\hbar^2}{2\mu}\nabla^2 \psi(\vec{r}) + V(r)\psi(\vec{r}) = E\psi(\vec{r}).
\]

对球对称势，角动量 $L^2$ 和 $L_z$ 守恒，可分离变量：
\[
\psi(\vec{r}) = R_l(r) Y_{lm}(\theta, \phi).
\]

代入球坐标下的拉普拉斯算符：
\[
\nabla^2 = \frac{1}{r^2}\frac{\partial}{\partial r}\left(r^2 \frac{\partial}{\partial r}\right) - \frac{L^2}{\hbar^2 r^2}, \qquad L^2 Y_{lm} = l(l+1)\hbar^2 Y_{lm}.
\]

径向部分变为：
\[
\frac{1}{r^2}\frac{d}{dr}\left(r^2 \frac{dR_l}{dr}\right) + \left[ \frac{2\mu}{\hbar^2}\bigl(E - V(r)\bigr) - \frac{l(l+1)}{r^2} \right] R_l = 0.
\]

令 $R_l(r) = u_l(r)/r$，计算：
\[
\frac{dR_l}{dr} = \frac{u_l'}{r} - \frac{u_l}{r^2}, \qquad \frac{d}{dr}\left(r^2 \frac{dR_l}{dr}\right) = r u_l''.
\]
代回，$1/r^2$ 的项与离心势垒的 $l(l+1)/r^2$ 相互抵消，奇迹般地得到一维形式的方程：
\[
\frac{d^2 u_l}{dr^2} + \left[ \frac{2\mu}{\hbar^2}\bigl(E - V(r)\bigr) - \frac{l(l+1)}{r^2} \right] u_l = 0.
\]
边界条件 $u_l(0) = 0$（否则 $R_l = u_l/r$ 在 $r = 0$ 发散）。

对 s 波（$l = 0$），离心势垒项消失，方程简化为：
\[
\frac{d^2 u}{dr^2} + \frac{2\mu}{\hbar^2}\bigl(E - V(r)\bigr) u = 0.
\]
这正是题目中的两个区间的方程来源。$r < a$ 时 $V = -3V_0$，$r > a$ 时 $V = 0$。
\end{derivationbox}

\begin{tipbox}[球 Bessel 恐惧消除器]
你担心分波法要背球 Bessel 函数 $j_l(x), n_l(x)$？对于 $l = 0$ 的 s 波，完全不需要：
\[
j_0(x) = \frac{\sin x}{x}, \qquad n_0(x) = -\frac{\cos x}{x}.
\]
代入 $R_0 = u_0/r$ 且 $u_0 \propto \sin(kr + \delta_0)$，你看到的只是最普通的正弦函数。

什么时候真需要球 Bessel？只有当题目要求你算 $l \ge 1$（p 波、d 波等）的相移，或者势的形式复杂（如 $V(r) \propto r^{-1}$ 库仑势要用合流超几何函数）。

考场策略：低能球方势阱散射，只考 $l = 0$。把 $u(r) = rR(r)$ 的方程记住，不用碰任何特殊函数。
\end{tipbox}

\begin{derivationbox}[s 波分波法求解]
在质心系中，约化质量 $\mu = m_n/2$。对单态通道，令 $\psi(r) = u(r)/r$，径向方程：
\[
\begin{cases}
\dfrac{d^2 u}{dr^2} + \alpha^2 u = 0, & r < a \\[6pt]
\dfrac{d^2 u}{dr^2} + k^2 u = 0, & r > a
\end{cases}
\qquad
\alpha = \sqrt{\frac{2\mu(E + 3V_0)}{\hbar^2}}, \quad k = \sqrt{\frac{2\mu E}{\hbar^2}}.
\]

边界条件 $u(0) = 0$，以及 $r \to \infty$ 处的散射边界条件给出：
\[
u_1(r) = A\sin(\alpha r), \qquad u_2(r) = B\sin(kr + \delta_0).
\]

在 $r = a$ 处匹配 $u$ 与 $u'$：
\begin{align*}
A\sin(\alpha a) &= B\sin(ka + \delta_0), \\
A\alpha\cos(\alpha a) &= Bk\cos(ka + \delta_0).
\end{align*}

两式相除：
\[
\tan(ka + \delta_0) = \frac{k}{\alpha}\tan(\alpha a) \quad \Longrightarrow \quad \delta_0 = \arctan\!\left[ \frac{k}{\alpha}\tan(\alpha a) \right] - ka.
\]
\end{derivationbox}

\begin{derivationbox}[相移的 $n\pi$ 不确定性与物理唯一性]
从 $\tan(ka + \delta_0) = C$ 只能确定 $ka + \delta_0$ 模 $\pi$，因此形式上 $\delta_0$ 可以加减任意 $n\pi$。但这完全不构成物理上的不唯一，原因有三：
\begin{enumerate}
    \item \textbf{可观测量不变：}散射振幅 $f(\theta) = \frac{1}{k}e^{i\delta_0}\sin\delta_0$ 在 $\delta_0 \to \delta_0 + n\pi$ 下，$e^{i(\delta_0+n\pi)}\sin(\delta_0 + n\pi) = (-1)^n e^{i\delta_0} \cdot (-1)^n \sin\delta_0 = e^{i\delta_0}\sin\delta_0$，严格不变。截面 $\sigma \propto |f|^2$ 更不变。
    \item \textbf{幺正性条件：}S 矩阵的幺正性要求 $S_l = e^{2i\delta_l}$，$\delta_l$ 只要保证 $S_l$ 唯一即可。$e^{2i(\delta_l + n\pi)} = e^{2i\delta_l}$，所以 $n\pi$ 加减是纯粹的``坐标规范''。
    \item \textbf{标准约定唯一化：}物理上要求势趋于零时相移趋于零，即 $\delta_l \xrightarrow{V \to 0} 0$。这选出了唯一的主值分支。更实操地说：让势从小逐渐增强，相移应该连续变化地从 0 开始；不能跳跃 $n\pi$。
\end{enumerate}

附：三角函数``模多少'' 速查卡——看到什么函数就模什么：

\begin{tabular}{lll}
\textbf{方程} & \textbf{解的周期性} & \textbf{你应该模} \\
$\tan x = C$ & $\tan(x + \pi) = \tan x$ & $\pi$ \\
$\sin x = C$ & $\sin(x + 2\pi) = \sin x$ & $2\pi$ \\
$\cos x = C$ & $\cos(x + 2\pi) = \cos x$ & $2\pi$ \\
$e^{ix} = C$ & $e^{i(x+2\pi)} = e^{ix}$ & $2\pi$ \\
\end{tabular}

考场实操：匹配条件得到的是 $\tan(ka + \delta_0) = \text{某数}$，所以 $ka + \delta_0$ 只能确定到模 $\pi$。但解 $\delta_0$ 时你写 $\arctan$ 就够了，因为物理结果只依赖于 $\sin^2\delta_0$。

一句话：$\delta_0$ 的 $n\pi$ 不确定性跟方位角 $\phi$ 的 $2\pi$ 周期性一样——坐标可以有歧义，但物理结果唯一。
\end{derivationbox}

\begin{derivationbox}[低能极限 $E \to 0$]
当 $E \to 0$ 时，$k \to 0$，$\alpha \to \alpha_0 \equiv \sqrt{6\mu V_0/\hbar^2}$。对 $\delta_0$ 做泰勒展开：
\[
\delta_0 \approx k\, \frac{\tan(\alpha_0 a)}{\alpha_0} - ka = ka \left[ \frac{\tan(\alpha_0 a)}{\alpha_0 a} - 1 \right].
\]

定义散射长度
\[
a_s \equiv -\lim_{k\to 0} \frac{\delta_0}{k} = a \left[ 1 - \frac{\tan(\alpha_0 a)}{\alpha_0 a} \right],
\]
则 $\delta_0 \approx -k a_s$。
\end{derivationbox}

\begin{insightbox}[散射长度的物理含义——为什么 Born 近似里没有它？]
物理图像：在极低能量下，入射粒子的德布罗意波长 $\lambda = 2\pi/k \to \infty$，远大于势的作用范围 $a$。粒子``感受'' 不到势的内部细节，只被一个等效的硬球边界散射。散射长度 $a_s$ 就是这个等效硬球的半径。

\textbf{为什么 Born 近似没有它？}
\begin{itemize}
    \item Born 近似是微扰展开：$f = f^{(1)} + f^{(2)} + \cdots$，假设势是小的扰动。
    \item 散射长度是非微扰的、普适的低能参数：无论势多强，只要 $k \to 0$，所有细节都被 $a_s$ 吸收。
    \item Born 近似在低能下失效，因为它不能正确描述 s 波相移的 $k \to 0$ 行为（Born 给出的 $f^{(1)}$ 通常有限，但真实的 $f \approx -a_s$ 与 $k$ 无关）。
\end{itemize}

定量关系：对 s 波，$f(\theta) = \frac{1}{k}e^{i\delta_0}\sin\delta_0 \approx \frac{\delta_0}{k} = -a_s$（当 $k \to 0$）。截面：
\[
\sigma_0 = 4\pi|f|^2 \xrightarrow{k\to 0} 4\pi a_s^2.
\]
这被称为有效范围理论（Effective Range Theory）的零阶结果。在冷原子物理中，散射长度是调控原子间相互作用的核心参数（Feshbach 共振就是通过磁场调节 $a_s$ 从 $-\infty$ 到 $+\infty$）。
\end{insightbox}

\begin{derivationbox}[对称化、非极化平均与总截面]
\subsubsection*{极化、非极化与部分极化——考场不害怕指南}

在涉及自旋的散射问题中，``极化'' 状态决定你对初态自旋做什么样的平均。核心只有三种情况：

\textbf{1. 非极化（Unpolarized）——完全混合态}

两个中子各自自旋方向完全随机，没有任何关联。用密度矩阵描述：
\[
\rho = \frac{1}{4} I_4 = \frac{1}{4}\bigl( |{\uparrow\uparrow}\rangle\langle{\uparrow\uparrow}| + |{\uparrow\downarrow}\rangle\langle{\uparrow\downarrow}| + |{\downarrow\uparrow}\rangle\langle{\downarrow\uparrow}| + |{\downarrow\downarrow}\rangle\langle{\downarrow\downarrow}| \bigr).
\]
这里 $I_4$ 是 $4 \times 4$ 单位矩阵，$1/4$ 是归一化（四个态等概率）。关键：$\rho$ 没有非对角元，意味着各态之间不相干。

转到耦合基 $|s, m_s\rangle$：单态 $|0, 0\rangle$（1 个），三重态 $|1, m\rangle$（3 个），共 4 个正交归一态。由于非极化对各方向一视同仁，这 4 个耦合基态也等概率：
\[
P(s = 1, \text{三重态}) = \frac{3}{4}, \qquad P(s = 0, \text{单态}) = \frac{1}{4}.
\]
考场操作：分别算各通道截面，按概率加权求和。

\textbf{2. 极化（Polarized）——纯态或特定取向}

例如``两中子均沿 $z$ 轴向上''，初态就是 $|{\uparrow\uparrow}\rangle = |1, 1\rangle$（纯三重态）。低能下三重态被统计禁戒，散射截面直接为零。

再如``一个向上、一个向下''，要仔细看是乘积基的 $|{\uparrow\downarrow}\rangle$ 还是耦合基的 $|1, 0\rangle$ 或 $|0, 0\rangle$。前者可写成 $\frac{1}{\sqrt{2}}(|1, 0\rangle + |0, 0\rangle)$，各以 $1/2$ 概率处于三重态或单态。

\textbf{3. 部分极化（Partially polarized）}

密度矩阵既有对角元（概率）又有非对角元（相干），一般形式：
\[
\rho = \sum_i p_i |\psi_i\rangle\langle\psi_i|, \qquad \sum_i p_i = 1.
\]
截面计算：$\sigma = \text{Tr}(\rho\, \hat{\sigma})$，其中 $\hat{\sigma}$ 是散射截面算符。

\textbf{题目说法速查：}

\begin{tabular}{lll}
\textbf{题目说法} & \textbf{初态密度矩阵} & \textbf{你该做的事} \\
``非极化'' & $\rho = I_4/4$ & 等概率平均四个自旋态 \\
``均沿 $z$ 向上'' & $\rho = |{\uparrow\uparrow}\rangle\langle{\uparrow\uparrow}|$ & 纯三重态，低能截面为零 \\
``部分极化'' & 一般 $\rho$（题目给定） & $\sigma = \text{Tr}(\rho\, \hat{\sigma})$ \\
\end{tabular}

对全同费米子的单态（空间对称），散射振幅需对称化：
\[
f(\theta) \longrightarrow f_S(\theta) = f(\theta) + f(\pi - \theta).
\]
物理含义：探测器在 $\theta$ 方向捕获一个中子，它可能是原来的``入射'' 中子直接散射到 $\theta$，也可能是原来的``靶'' 中子被反冲到 $\theta$（即入射中子去了 $\pi - \theta$）。由于两中子全同，这两种情况不可区分，必须将两种路径的振幅叠加。

s 波振幅与角度无关，$f(\theta) = \frac{1}{k}e^{i\delta_0}\sin\delta_0$，故
\[
f_S(\theta) = \frac{2}{k}\, e^{i\delta_0}\sin\delta_0.
\]
单态微分截面：
\[
\sigma_S(\theta) = |f_S(\theta)|^2 = \frac{4}{k^2}\sin^2\delta_0.
\]

\textbf{非极化平均——逐步写清楚：}

初态是非极化的，意味着四个自旋基态各以 $1/4$ 概率出现：

\begin{tabular}{llll}
\textbf{乘积基} & \textbf{耦合基} & \textbf{总自旋} & \textbf{概率} \\
$|{\uparrow\uparrow}\rangle$ & $|1, 1\rangle$ & $s = 1$（三重态） & $1/4$ \\
$|{\uparrow\downarrow}\rangle$ & $\frac{1}{\sqrt{2}}(|1, 0\rangle + |0, 0\rangle)$ & 混合 & $1/4$ \\
$|{\downarrow\uparrow}\rangle$ & $\frac{1}{\sqrt{2}}(|1, 0\rangle - |0, 0\rangle)$ & 混合 & $1/4$ \\
$|{\downarrow\downarrow}\rangle$ & $|1, -1\rangle$ & $s = 1$（三重态） & $1/4$ \\
\end{tabular}
\end{derivationbox}
